\section{Introduction and Problem Definition}
\label{ch:intro}

\subsection{Project Background}
After an earthquake, fire or collapse the mobile network and internet link are usually the first to fail or saturate, yet connectivity survives in pockets (a person near a window with one bar of signal, an access point on a UPS, a rescue team's kit) that help only their owner. \textbf{EcoMesh} turns every participating device into both a sender and a relay, so a message can travel from device to device until it reaches its addressee, a \emph{gateway} that still has an outside link, or a \emph{rescue} device, without the internet and without a server.

\subsection{Problem Statement}
WhatsApp, Signal, Telegram and SMS all need a distant server or tower; when a Wi-Fi router loses its uplink they all stop, although the phones on it can still exchange packets. Trapped people then cannot report their condition, share the one weak signal, coordinate or be located quickly. Existing offline systems either flood messages or route by hop count, give no priority to emergencies, publish few measurements, and in one widely deployed case were shown to lack effective confidentiality and authenticity~\cite{albrecht2021bridgefy} (Chapter~\ref{ch:lit}).

\subsection{Motivation}
Survival depends on how soon rescuers learn who is inside, where, and in what condition. Sharing one person's connectivity with everyone nearby and pushing an SOS ahead of routine traffic at every hop shortens that delay for trapped occupants, rescue teams, and the hospitals, campuses, hostels, factories and stadiums that could install it in advance.

\subsection{Objectives}
O1 \emph{Self-organising mesh}: neighbours are found automatically and every node learns the whole graph within 15\,s. O2 \emph{Weighted routing} that outperforms hop count on identical scenarios. O3 \emph{Priority}: an SOS overtakes routine traffic at every hop, with a median latency at least 3$\times$ lower under load. O4 \emph{Reliability}: at least 90\,\% delivery on 40 nodes at 20\,\% loss, and a new route within 10\,s of a relay failure. O5 \emph{Store-and-forward}: a message without a route is kept, never dropped, and is delivered when a route returns. O6 \emph{Security}: payloads stay confidential and tamper-evident, and a forged SOS is rejected. O7 \emph{Evidence}: every metric is computed from the nodes' own logs. O8 \emph{Integration}: a web front end, a REST and WebSocket back end and the engine running together on several devices.

\subsection{Scope of the Project}
\textbf{Included:} discovery, a binary wire protocol, link-state gossip, Dijkstra routing over a link-quality cost, per-hop priority, TTL, duplicate suppression, acknowledgement with retry, heartbeats, node roles, bounded store-and-forward, AES-256-GCM and Ed25519, a SQLite log per node, \texttt{asyncio} concurrency, a simulator with metrics, and survivor and rescue screens. \textbf{Excluded:} a real Bluetooth LE or Wi-Fi Direct transport, NAT traversal, voice calls, video and file transfer, machine learning and a native mobile app. \textbf{Users:} the trapped person, the relay user, the gateway operator and the rescue coordinator. \textbf{Environment and scale:} laptops running Python~3.10 on a shared network, five live devices and up to 65 simulated nodes. One design decision is stated plainly: the mesh runs over LAN sockets that stand in for Bluetooth LE or Wi-Fi Direct, and each node may talk only to the neighbours listed for it, which imitates limited radio range. Discovery, routing, reliability and encryption sit above the transport, so a radio transport can replace the sockets without changing them.

\subsection{Proposed Solution}
Each device runs one node, a Python process whose \texttt{asyncio} loop handles discovery, gossip, receiving, sending, store-and-forward and a small web server. A message leaves by one of three exits: a \textbf{gateway} with an outside link (reaching any gateway counts as delivery); a \textbf{rescue} device that comes into range and receives every pending SOS; or \textbf{store-and-forward}, where a node with no route keeps the message, so a person walking out carries others' messages out. Paths are chosen by Dijkstra's algorithm on a graph whose edge weights penalise weak and failing links (Chapter~\ref{ch:alg}).

\subsection{Novelty and Originality}
EcoMesh combines \textbf{priority-aware, link-quality-weighted routing with a delay-tolerant fallback, and evaluates it by measurement}. SOS priority is applied again at every hop. The link cost is derived from measured HELLO reception and from acknowledgement failures, and a failure raises the cost immediately. A single transport interface lets the same engine run on a simulator, on loopback and on separate laptops. Every reported number is a query over the nodes' logs, and weighted and hop-count routing are compared on the same seeds.

\subsection{Feasibility Analysis}
\textbf{Technical:} all components are mature open-source software, and the largest risk, debugging routing over live sockets, was removed by building the simulator first. \textbf{Economic:} the system uses existing laptops and free software and needs no server, so there is no running cost. \textbf{Operational:} one command starts a node, and the survivor screen has one prominent control. \textbf{Time and resources:} three members over one semester, about 336 person-hours, delivered in four tagged phases.

\subsection{Social Relevance and Practical Applications}
The realistic path is \emph{installation before an incident}: an institution puts EcoMesh on staff and resident devices as part of its safety plan and rescue teams carry it as standard kit, so the rescue exit exists even if no civilian has it. Survivors are found faster at no infrastructure cost.

\subsection{Success Criteria and Key Performance Indicators}
\label{sec:kpi}
Table~\ref{tab:kpi} lists the quantified requirements of the SRS, fixed before testing began; Section~\ref{sec:kpival} checks each against measured data.
\begin{table}[!htbp]
\caption{Key performance indicators}
\label{tab:kpi}
\footnotesize
\begin{tabularx}{\textwidth}{@{}>{\raggedright\arraybackslash}p{3.6cm}>{\raggedright\arraybackslash}X>{\raggedright\arraybackslash}p{1.9cm}>{\raggedright\arraybackslash}p{3.5cm}@{}}
\toprule
\textbf{KPI} & \textbf{Definition} & \textbf{Target} & \textbf{Measured by}\\ \midrule
K1 Delivery & Delivered $\div$ sent, 40 nodes, 20\,\% loss & $\ge$ 90\,\% & seeded simulator runs\\
K2 Convergence & Every node's graph holds every node & $\le$ 15\,s & simulator, live\\
K3 Recovery & Relay death until a route avoiding it & $\le$ 10\,s & 50 trials, live\\
K4 SOS advantage & Median routine $\div$ median SOS latency under load & $\ge$ 3$\times$ & burst scenario\\
K5 Store-and-forward & Retained messages flushed on return & 100\,\% & simulator, live\\
K6 Weighted vs hop count & Delivery on identical seeds & weighted better & paired runs\\
K7 Scale & Nodes within the targets & $\ge$ 50 & sweep 10--65 nodes\\
K8 UI freshness & Event to screen & $\le$ 2\,s & API time, poll interval\\
K9 Correctness & Tests passed & 100\,\% & repeated runs\\
K10 Security & No relay decryption; forged SOS and tampering rejected & all hold & automated tests\\
\bottomrule
\end{tabularx}
\end{table}
